
\documentclass{article}
\usepackage{colortbl}
\usepackage{makecell}
\usepackage{multirow}
\usepackage{supertabular}

\begin{document}

\newcounter{utterance}

\centering \large Interaction Transcript for game `cladder', experiment `full\_v1.5\_default', episode 7019 with reflexion{-}Qwen3.5{-}9B.
\vspace{24pt}

{ \footnotesize  \setcounter{utterance}{1}
\setlength{\tabcolsep}{0pt}
\begin{supertabular}{c@{$\;$}|p{.15\linewidth}@{}p{.15\linewidth}p{.15\linewidth}p{.15\linewidth}p{.15\linewidth}p{.15\linewidth}}
    \# & \multicolumn{2}{c}{Player} && \multicolumn{2}{c}{Game Master} \\
    \hline

    \theutterance \stepcounter{utterance}  
    & & & \multicolumn{4}{p{0.6\linewidth}}{
        \cellcolor[rgb]{0.9,0.9,0.9}{
            \makecell[{{p{\linewidth}}}]{
                \texttt{\tiny{[P1$\langle$GM]}}
                \texttt{You have attempted this same task before and did not succeed. These are the plans you wrote after those attempts.} \\
\texttt{Trial \#1: New plan: I failed because I provided an incomplete answer that cut off mid{-}sentence, preventing the system from registering a valid "Yes" or "No" conclusion. In causal inference problems involving confounders (like \$Hwax\$) and multiple paths to an outcome (\$lirg\$), the marginal correlation (Method 2) is biased by the confounder, whereas stratification (Method 1) blocks the backdoor path through \$gyzp\$, isolating the direct effect of \$jyka\$. Next time, I will immediately start with a definitive "No" (since Method 2 is biased), followed by a complete, concise explanation that explicitly identifies \$Hwax\$ as a confounder and explains how conditioning on \$gyzp\$ removes the bias, ensuring the response is fully formed before submission.} \\
\\ 
\\ 
\texttt{You are an expert in causal inference. The following question is not a typical commonsense query, but rather a meticulously designed question created by a professor specializing in causal inference, intended to assess the students' mastery of the course content.} \\
\texttt{Imagine a self{-}contained, hypothetical world with only the following conditions, and without any unmentioned factors or causal relationships: Hwax has a direct effect on gyzp and jyka. Jyka has a direct effect on lirg. Gyzp has a direct effect on lirg. Method 1: We look at how jyka correlates with lirg case by case according to gyzp. Method 2: We look directly at how jyka correlates with lirg in general. To understand how jyka affects lirg, is it more correct to use the Method 1 than Method 2?} \\
\texttt{Start your answer with 'Yes' or 'No', followed by additional reasoning or evidence to support your explanation.} \\
            }
        }
    }
     \\ \\

    \theutterance \stepcounter{utterance}  
    & \multicolumn{4}{p{0.6\linewidth}}{
        \cellcolor[rgb]{0.561,0.737,0.561}{
            \makecell[{{p{\linewidth}}}]{
                \texttt{\tiny{[P1$\rangle$GM]}}
                \texttt{Yes} \\
\\ 
\texttt{In this hypothetical world, \$Hwax\$ acts as a confounder because it has direct causal paths to both the treatment (\$jyka\$) and the outcome (\$lirg\$), creating a backdoor path: \$jyka \textbackslash{}leftarrow Hwax \textbackslash{}rightarrow gyzp \textbackslash{}rightarrow lirg\$ (and potentially \$jyka \textbackslash{}leftarrow Hwax \textbackslash{}rightarrow lirg\$ if \$Hwax\$ affects \$lirg\$ directly, though the prompt specifies \$Hwax \textbackslash{}rightarrow gyzp\$ and \$gyzp \textbackslash{}rightarrow lirg\$, the key is that \$Hwax\$ influences the distribution of \$gyzp\$ which influences \$lirg\$, and \$Hwax\$ also influences \$jyka\$). More critically, since \$Hwax\$ affects \$jyka\$ and \$gyzp\$ (which affects \$lirg\$), there is a non{-}causal association between \$jyka\$ and \$lirg\$ induced by \$Hwax\$.} \\
\\ 
\texttt{Method 2 calculates the marginal correlation between \$jyka\$ and \$lirg\$. Because \$Hwax\$ is a common cause of the variables in the causal chain leading to \$lirg\$ (via \$gyzp\$) and \$jyka\$, the marginal correlation will be biased by the confounding effect of \$Hwax\$. It conflates the true causal effect of \$jyka\$ on \$lirg\$ with the spurious correlation induced by \$Hwax\$.} \\
\\ 
\texttt{Method 1 involves stratifying by \$gyzp\$ (looking at the correlation case by case according to \$gyzp\$). By conditioning on \$gyzp\$, we block the path from \$Hwax\$ to \$lirg\$ that flows through \$gyzp\$ (\$Hwax \textbackslash{}rightarrow gyzp \textbackslash{}rightarrow lirg\$). Since \$gyzp\$ is a mediator on the path from \$Hwax\$ to \$lirg\$, and \$Hwax\$ is the source of the confounding bias relative to the \$jyka \textbackslash{}rightarrow lirg\$ relationship, conditioning on \$gyzp\$ isolates the variation in \$lirg\$ that is not explained by the \$Hwax \textbackslash{}rightarrow gyzp \textbackslash{}rightarrow lirg\$ pathway. Specifically, within strata of fixed \$gyzp\$, the effect of \$Hwax\$ on \$lirg\$ is held constant (blocked), allowing us to observe the direct effect of \$jyka\$ on \$lirg\$ without the bias introduced by the varying levels of \$Hwax\$ across the population. Therefore, Method 1 provides an unbiased estimate of the causal effect, while Method 2 does not.} \\
            }
        }
    }
    & & \\ \\

    \theutterance \stepcounter{utterance}  
    & & & \multicolumn{2}{p{0.3\linewidth}}{
        \cellcolor[rgb]{0.9,0.9,0.9}{
            \makecell[{{p{\linewidth}}}]{
                \texttt{\tiny{[GM$|$GM]}}
                \texttt{Yes} \\
            }
        }
    }
    & & \\ \\

    \theutterance \stepcounter{utterance}  
    & & & \multicolumn{2}{p{0.3\linewidth}}{
        \cellcolor[rgb]{0.9,0.9,0.9}{
            \makecell[{{p{\linewidth}}}]{
                \texttt{\tiny{[GM$|$GM]}}
                \texttt{yes} \\
            }
        }
    }
    & & \\ \\

    \theutterance \stepcounter{utterance}  
    & & & \multicolumn{2}{p{0.3\linewidth}}{
        \cellcolor[rgb]{0.9,0.9,0.9}{
            \makecell[{{p{\linewidth}}}]{
                \texttt{\tiny{[GM$|$GM]}}
                \texttt{game\_result = WIN} \\
            }
        }
    }
    & & \\ \\

\end{supertabular}
}

\end{document}
